\section{Frozen computational experiment}
\label{sec:experiment}

\subsection{Source lock and splits}

The protocol was serialized before the sealed periods were generated.  Its
SHA-256 digest is
\begin{center}
\small\texttt{0c284a1b3610a3d772aa00c6a8b33161a8bc6814957a9968d5c80fb618eec399}.
\end{center}
The object is \(H_6|_{\Lam}\); the clock, scale, and offset are fixed as
\(T_p=\log\abs{\Lambda_{u,p}}\), scale one, offset zero.  The complex
rectangle is
\[
 \cR=\{s\in\C:-0.25\le\Re s\le0.30,\ \abs{\Im s}\le20\}.
\]
Periods 1--8 form development, 9--12 validation, 13--16 the sealed test, and
17--20 post-test robustness.  The matching tolerance is \(0.02\).  The
training family contains cutoff-8 zeros that match cutoff 7 within this
tolerance and lie at least \(0.01\) from the contour.

Prime tables, primality labels, Riemann-zero tables, evaluations of zeta or xi,
target spectral fitting, and post-validation rescaling were forbidden.  This
firewall makes the experiment a self-consistency test of one dynamical
candidate, not a goodness-of-fit exercise.

\subsection{Complete orbit ledger}

Every primitive closed \(A\)-word is enumerated exactly up to cyclic rotation.
Equation \eqref{eq:sqrt-recurrence} then lifts it at 80 decimal digits.  For
each orbit, the ledger stores its canonical state and sign words, coordinates,
recurrence residual, a posteriori contraction bound, monodromy, unstable
multiplier, orientation, instability length, and generating-function action.
The catalogue is complete within the symbolic survivor at every stated
cutoff.

The two neighboring catalogues \(a=5.9,6.1\) reuse the same words as numerical
continuations only.  On the common boxes, their contraction diagnostic uses
\[
 \kappa(a)=\frac{2}{\sqrt{17}}\sqrt{\frac6a}.
\]
No common-survivor theorem is inferred from their convergence.

\subsection{Two coefficient constructions}

Primitive-factor multiplication gives the coefficients of
\eqref{eq:section} directly.  An independent trace recursion uses
\begin{equation}
 F_j^{(\kappa)}(s)=
 \sum_{\substack{p\in\cP\\n_p\mid j}}
 n_p\,\sigma_p^{\kappa j/n_p}\e^{-s(j/n_p)T_p},
 \label{eq:traces}
\end{equation}
followed by
\begin{equation}
 d_0(s)=1,\qquad
 d_n(s)=-\frac1n\sum_{j=1}^nF_j^{(\kappa)}(s)d_{n-j}(s),
 \qquad
 D_{\kappa,N}(s)=\sum_{n=0}^N d_n(s).
 \label{eq:cumulant}
\end{equation}
Both constructions use the same complete orbit ledger but independent
algebraic paths.  Agreement is tested at fixed complex probes and at every
reported zero.

\subsection{Zeros and numerical winding audit}

Float64 discovery combines grid-cell phase windings and local modulus minima,
then applies bounded Newton refinement.  Each located zero is refined at 80
decimal digits and checked with both coefficient constructions.  The sampled
argument change around \(\partial\cR\) is evaluated at 4,096, 8,192, and
16,384 contour points.  We record agreement across resolutions, the minimum
sampled boundary modulus, and the maximum phase step.

These winding counts are numerical diagnostics, not interval certificates.
Even when all three grids agree and every sampled phase step is below \(\pi\),
an analytic derivative bound or interval enclosure would be required to rule
out unresolved behavior between samples.

\subsection{Controls}

The frozen control panel contains global period and length shuffles,
same-density random lengths, positive random factor weights, random phases, an
exact constant-roof parent, and numerical continuations at \(a=5.9,6.1\).
For the exact parent, every step has length
\(\ell_*=\log L_1\simeq1.96734662909421\), the negative fixed-orbit length,
so \(T_p=n_p\ell_*\) and its determinant is
\(\det(I-z\e^{-s\ell_*}A)\); coefficients above degree four vanish exactly.
Random controls use seeds 20260805--20260807.  Root retention is measured from
cutoff 8 to 16.  A second statistic is the degree-9--16 coefficient
\(\ell^1\) tail at the common H\'enon cutoff-16 positive zero.  Since the
randomized orbit data need not define H\"older potentials, this panel diagnoses
non-generic cancellation but cannot alone identify its mechanism.
